\documentclass{exam-zh}

\usepackage{graphicx}

\title{2026 年普通高等学校招生全国统一考试}
\subject{数学}

\everymath{\displaystyle}

\begin{document}

\secret

\maketitle

\begin{notice}
\item 答卷前, 考生务必将自己的姓名、考生号、考场号、座位号填写在答题卡上.
\item 回答选择题时, 选出每小题答案后, 用铅笔把答题卡上对应题目的答案标号涂黑,
  如需改动, 用橡皮擦干净后, 再选涂其他答案标号.
  回答非选择题时, 将答案写在答题卡上. 写在本试卷上无效.
\item 考试结束后, 将本试卷和答题卡一并交回.
\end{notice}

\section{选择题：本题共 8 小题, 每小题 5 分, 共 40 分, 在每小题给出的四个选项中, 只有一项是符合题目要求的. }
\begin{question} %1
  样本数据 6, 8, 4, 5, 12 的中位数为 \paren
  \begin{choices}
  \item 5
  \item 6
  \item 8
  \item 9
  \end{choices}
\end{question}

\begin{question} %2
  已知平面向量 $\vec{a}, \vec{b}$ 不共线, 且 $2 \vec{a} + y \vec{b} = x \vec{a}
  - 3 \vec{b}$, 则 \paren
  \begin{choices}
  \item $x=2, y=-3$
  \item $x=-2, y=3$
  \item $x=2, y=3$
  \item $x=-2, y=-3$
  \end{choices}
\end{question}

\begin{question} %3
  已知集合
  $A=\left\{\sin\frac{7\pi}{6},\cos\frac{5\pi}{3},\tan\frac{5\pi}{4}\right\}$,
  集合 $B=\left\{-\frac{\sqrt{3}}{2},-\frac{1}{2},1\right\}$, 则 $A\cap B=$ \paren
  \begin{choices}
  \item $\left\{-\frac{\sqrt{3}}{2},-\frac{1}{2}\right\}$
  \item $\left\{-\frac{\sqrt{3}}{2}, 1\right\}$
  \item $\left\{-\frac{1}{2}, 1\right\}$
  \item $\left\{-\frac{\sqrt{3}}{2},-\frac{1}{2}, 1\right\}$
  \end{choices}
\end{question}

\begin{question}%4
  曲线 $y=5x+8\ln x$ 在点 $(1,5)$ 处的切线方程为 \paren
  \begin{choices}
  \item $y=3x+2$
  \item $y=5x$
  \item $y=8x-3$
  \item $y=13x-8$
  \end{choices}
\end{question}

\begin{question}%5
  已知抛物线 $C_{1}: y^{2}=2 p_{1} x\left(p_{1}>0\right)$ 和
  $C_{2}:x^2=2 p_{2} y\left(p_{2}>0\right)$
  均经过点 $(4,8)$, 则 $C_{1}$ 的焦点与 $C_{2}$ 的焦点之间的距离为 \paren
  \begin{choices}
  \item $12$
  \item $4\sqrt{5}$
  \item $6$
  \item $\frac{\sqrt{65}}{2}$
  \end{choices}
\end{question}

\begin{question}%6
  已知函数 $f(x)=\frac{x+2}{e^{x}+a}$ 的最大值为 1, 则 $a=$\paren
  \begin{choices}
  \item $\frac{1}{2} $
  \item $1$
  \item $\frac{3}{2} $
  \item $2$
  \end{choices}
\end{question}

\begin{question}%7
  一百零八塔位于宁夏回族自治区青网峡市, 以其独特的建筑格局和深远的历史文化间名遇连,
  该塔群共有 108 座塔, 依山势自上而下排成 12 行,
  设第$i$行中塔的座数记为 $a_{1}(l=1,2,\cdots,12)$,
  其中 $a_{1}=1, a_{2}=a_{3}=3,a_{4}=a_{5}=5$,
  且 $a_{6},a_{7},\cdots,a_{12}$ 是一个首项为 7、公差为 2 的等差数列.
  将 $a_{1}, a_{2}, \cdots, a_{12}$ 分为 6 组,每 组 2 个数,
  使得每组的 2 个数之和可构成一个项数为 6 且公差为 $d(d>0)$ 的等差数列, 则 $d=$\paren
  \begin{choices}
  \item 2
  \item 4
  \item 6
  \item 8
  \end{choices}
\end{question}

\begin{question}%8
  设 $U=\left\{\left(x_{1},x_{2},x_{3}\right) \mid x_{i}
  \in\{-2,-1,1,2\}, i=1,2,3\right\}$
  为空间中 64 个点构成的集合, 点 $P(1,1,1)$, 记样本空间 $\Omega=\complement_{U}(P)
  $. 从 $\Omega$ 中随机取一个点.
  定义随机变量 $X$ 如下：对 $\Omega$ 中的每个点 $A\left(x_{1},x_{2},x_{3}\right)$,
  令 $X(A)=x_{1}+x_{2}+x_{3}$, 则 $X$ 的数学期望为 \paren
  \begin{choices}
  \item $-\frac{1}{21}$
  \item $-\frac{1}{63}$
  \item $0$
  \item $\frac{1}{7}$
  \end{choices}
\end{question}

\section{选择题：本题共 3 小题, 每小题 6 分, 共 18 分.
  在每小题给出的四个选项中有多项符合题目要求. 全部选对的得 6 分,
部分选对的得部分分、有选错的得 0 分.}

\begin{question}%9
  设 $z=3+2\iu$, 则 \paren
  \begin{choices}
  \item $\overline{z}=3-2\iu$
  \item $|z|=5$
  \item $z^{2}=5+12\iu$
  \item $\frac{z+3}{z-1}\in\mathbb{R}$
  \end{choices}
\end{question}

\begin{question}%10
  在空间中, $A$、$B$为两个定点, 动点 $C$ 到直线 $AB$ 的距离为 2, 动点 $D$ 到直线 $AB$ 的距离为 1.
  若二面角 $C-AB-D$ 为 $60^{\circ}$, 则 \paren
  \begin{choices}
  \item $\angle CAD\geqslant60^{\circ}$
  \item $CD\geqslant\sqrt{3}$
  \item 当 $AB\perp CD$ 时, $CD\perp$ 平面 $ABD$
  \item 当 $AB\perp$ 平面 $ACD$ 时, $AC\perp AD$
  \end{choices}
\end{question}

\begin{question}
  已知圆 $C_{1}:(x+1)^{2}+y^{2}=1$, 圆 $C_{2}:(x-1)^{2}+y^{2}=1$, 圆
  $C_{3}:x^{2}+(y-\sqrt{3})^{2}=1$,
  直线 $l:y=kx+b$ 与 $C_{1},C_{2},C_{3}$ 均有两个交点.
  记 $l$截$C_{1},C_{2},C_{3}$截得的弦长分别为 $s_{1}, s_{2}, s_{3}$, 则 \paren
  \begin{choices}
  \item $k$ 可以取任意实数
  \item 满足 $s_{1}=s_{2}=s_{3}$ 的直线 $l$ 共有 3 条
  \item 满足 $s_{1}+s_{2}+s_{3}=3$ 的直线 $l$ 多于 3 条
  \item 当 $b=0$ 时, $s_{1}+s_{2}+s_{3}$ 的最大值为 $\frac{2\sqrt{21}}{3}$
  \end{choices}
\end{question}

\section{填空题：本题共 3 小题, 每小题 5 分, 共 15 分. }

\begin{question}%12
  双曲线 $5x^{2}-6y^{2}=1$ 的离心率为 \fillin
\end{question}

\begin{question}%13
  已知 $f(x)=2\sin(ax+\theta)(a\in\mathbb{Z},0\leqslant\theta\leqslant2\pi)$ 是偶函数;
  $f(x)$ 在区间 $\left(0,\frac{\pi}{2}\right)$ 单调递增. 则 $\theta=$\fillin,
  $f\left(\frac{2\pi}{3}\right)=$\fillin.
\end{question}

\begin{question}
  设实数 $q$ 满足：存在数列 $\{a_{n}\}$, 使得对于任意 $n\in\mathbf{N}^{*}$,
  均有 $a_{1}+a_{2}+\cdots+a_{3n}=n^{2}+n$,
  且 $\{a_{n}\}$ 中有某连续 9 项 $a_{i},a_{i+1},\cdots,
  a_{i+8}$ 是公比为 $q$ 的等比数列, 则 $q$ 的最大值为 \fillin.
\end{question}

\newpage
\section{解答题：本题共 5 小题, 共 77 分. 解答应写出文字说明、证明过程或演算步骤. }

\begin{problem}[points=13]%15
  如图, 在直三棱柱 $ABC-A_{1}B_{1}C_{1}$ 中, $\angle ACB=90^{\circ}$,
  $AC=BC$, $D$, $E$ 分别为 $AB, AC$ 的中点.

  \begin{figure}
    \includegraphics*[scale=1, keepaspectratio]{figure/figure.png}
  \end{figure}

  \begin{enumerate}
    \item 证明：$DE\parallel$ 平面 $BCC_{1}B_{1}$
    \item 设 $CC_{1}=2$, 直线 $DE$ 与平面 $ACC_{1}A_{1}$ 所成的角为
      $45^{\circ}$, 求直线 $DE$ 到平面 $BCC_{1}B_{1}$ 的距离.
  \end{enumerate}
\end{problem}

\begin{problem}[points=15]%16
  已知在 $ \triangle A B C $ 中, $ A B=3, B C=2 \sqrt{3}, \cos
  B=\frac{\sqrt{3}}{3} $
  \begin{enumerate}
    \item 求 $\cos A$;
    \item 设 $D,E$ 两点满足：$D$ 在 $BA$ 的延长线上, $DE\parallel BC$,
      $AE\perp AC$.
      若 $DE=\sqrt{6}$, 求 $CE$.
  \end{enumerate}
\end{problem}

\begin{problem}[points=15]
  设整数  $N \geqslant 2 $ .某同学用一个球进行投篮练习, 至多投篮  $N$  次,
  当且仅当投中 1 次时或 $ N $ 次均末投中时, 停止练习.
  设该同学每次投中的概率为 $ p(0<p<1) $ , 各次投中与否相互独立.
  记 $ X $ 为停止练习时该同学的投篮次数.

  \begin{enumerate}
    \item 当 $N=4, p=\frac{1}{3}$ 时, 求  $X$ 的分布列：
    \item 设 $k, m$ 均为自然数.
      \begin{enumerate}
        \item 当 $ k \leqslant N-1$  时, 求 $P(X>k)$：
        \item 当 $ k+m \leqslant N-1 $ 时, 证明：$P(X>k+m \mid X>k)=P(X>m)$.
      \end{enumerate}
  \end{enumerate}
\end{problem}

\begin{problem}[points=17]
  己知椭圆 $ C: \frac{x^{2}}{a^{2}}+\frac{y^{2}}{b^{2}}=1(a>b>0) $ 的左焦点为 $F(-1,0)$ ,
  离心率为  $\frac{1}{2}$.

  \begin{enumerate}
    \item 求 $C$ 的方程：
    \item 设 $O$ 为坐标原点, 过 $F$ 且斜率大于 0 的动直线 $l$ 与 $C$ 交于 $P, Q$两点,
      其中 $ Q$ 在第三象限, 直线 $PO$  与 $C$  的另一个交点为 $R$.

      \begin{enumerate}
        \item 若 $\triangle PQR$  的面积是 $\triangle PFO$ 的面积的 3
          倍, 求  $l$  的方程;
        \item 求  $\tan \angle PQR$ 的最小值.
      \end{enumerate}

  \end{enumerate}
\end{problem}

\begin{problem}[points=17]
  已知函数 $f(x)$  的定义域为 $\mathbb{R} $ , 且当 $x<0 $ 时,  $f(x)=2^{x} $ .
  对任意 $x_{0} \in \mathbb{R}$, 定义集合

  \begin{equation*}
    D\left(x_{0}\right)=\left\{d \in \mathbb{R} \mid
    f\left(x_{0}+d\right)>f\left(x_{0}\right)\right\}
  \end{equation*}

  \begin{enumerate}
    \item 若当  $x \geqslant 0$  时, $ f(x)=1-x$  , 求  $D(-1)$;
    \item 若  $f(x) $ 是奇函数, $ f\left(x_{1}\right) \leqslant
      f\left(x_{2}\right) $ ,
      且 $ x_{1} x_{2} \neq 0$, 证明：$D\left(x_{2}\right) \subseteq
      D\left(x_{1}\right)$：
    \item 设 $ f(x)$  满足：
      （1）若 $ f\left(x_{1}\right) \leqslant f\left(x_{2}\right)$, 则
      $D\left(x_{2}\right) \subseteq D\left(x_{1}\right)$;

      （2）当 $ 0<x<1$  时,  $ f(x)<f(0) $.

      \begin{enumerate}
        \item  证明：$ f(0) \geqslant 1 $;
        \item 证明：$ f(x)  $在区间$  (0,+\infty) $ 单调递增.
      \end{enumerate}
  \end{enumerate}
\end{problem}

\end{document}
